% !TeX program = xelatex
%
% Лекция-сопровождение к лабораторной работе № 4 «Генерация текста».
%
% Сборка:  xelatex mlg_lab4_generation.tex   (дважды)
% Оформление и макросы — в mlg_preamble.tex, общем для всех лекций курса.
%
% Числа и примеры текста получены на символьной n-граммной модели из
% playground'а «Языковая модель и температура» (корпус курса, 160 сегментов)
% и закреплены тестами репозитория gd_learn. Сам ноутбук работы —
% демонстрация TensorFlow на английском корпусе Шекспира; его собственных
% чисел на слайдах нет намеренно, см. слайд «Что делает ноутбук работы».
%
\documentclass[aspectratio=169,12pt]{beamer}

% Общая преамбула лекций курса. Подключается после \documentclass:
%
%     \documentclass[aspectratio=169,12pt]{beamer}
%     % Общая преамбула лекций курса. Подключается после \documentclass:
%
%     \documentclass[aspectratio=169,12pt]{beamer}
%     % Общая преамбула лекций курса. Подключается после \documentclass:
%
%     \documentclass[aspectratio=169,12pt]{beamer}
%     \input{mlg_preamble}
%
% Здесь только оформление и макросы — ни \title, ни текста слайдов, чтобы
% пять презентаций курса выглядели одинаково и правились в одном месте.
%
\usepackage{fontspec}
\usepackage{polyglossia}
\setmainlanguage{russian}
\setotherlanguage{english}
\setmainfont{Calibri}
\setsansfont{Calibri}
\setmonofont{Consolas}[Scale=0.92]

\usepackage{graphicx}
\usepackage{booktabs}
\usepackage{array}

% ---------------------------------------------------------------- оформление
\definecolor{ink}{HTML}{1B2433}
\definecolor{navy}{HTML}{2B4C6F}
\definecolor{brick}{HTML}{B03A2B}
\definecolor{moss}{HTML}{3E7D5A}
\definecolor{plum}{HTML}{6B5B95}
\definecolor{ash}{HTML}{8B9299}
\definecolor{paper}{HTML}{ECEFF3}

\usetheme{default}
\usecolortheme{default}
\setbeamercolor{structure}{fg=navy}
\setbeamercolor{frametitle}{fg=navy}
\setbeamercolor{normal text}{fg=ink}
\setbeamercolor{itemize item}{fg=navy}
\setbeamercolor{itemize subitem}{fg=ash}
\setbeamerfont{frametitle}{size=\Large,series=\bfseries}
\setbeamerfont{framesubtitle}{size=\small,series=\mdseries}
\setbeamercolor{framesubtitle}{fg=ash}
\setbeamertemplate{navigation symbols}{}
\setbeamertemplate{itemize item}{\raisebox{0.15ex}{\tiny$\bullet$}}
\setbeamertemplate{itemize subitem}{\raisebox{0.15ex}{\tiny--}}
\setbeamersize{text margin left=1.1cm, text margin right=1.1cm}

\setbeamertemplate{footline}{%
  \hbox{\begin{beamercolorbox}[wd=\paperwidth,ht=2.6ex,dp=1.4ex,leftskip=1.1cm,%
        rightskip=1.1cm]{}%
    {\color{ash}\scriptsize Основы машинного обучения \quad·\quad
     Цифровая лингвистика и локализация}\hfill
    {\color{ash}\scriptsize\insertframenumber}%
  \end{beamercolorbox}}}

% Врезка: связь слайда с лабораторной или с работой переводчика.
% Ширина считается точно: colorbox добавляет \fboxsep с каждой стороны,
% поэтому parbox должен быть на 2\fboxsep уже текстового поля — иначе
% на каждом слайде вылезает overfull hbox.
\newcommand{\врезка}[1]{%
  \par\vspace{0.5ex}\noindent
  \colorbox{paper}{%
    \parbox{\dimexpr\textwidth-2\fboxsep\relax}{%
      \vspace{0.7ex}\small #1\vspace{0.7ex}}}\par}

% Колонка таблицы заданной доли ширины, без растягивания пробелов.
% Объявлять её нужно здесь: внутри frame beamer пересканирует тело кадра,
% и #1 в определении вызывает «Illegal parameter number».
\newcolumntype{Л}[1]{>{\raggedright\arraybackslash}p{#1\textwidth}}

% Ссылка на интерактивную страницу к слайду. Playground'ы считают те же
% числа на том же корпусе, что и лабораторные, поэтому на них можно
% ссылаться прямо с лекции, не оговаривая расхождений.
\hypersetup{colorlinks=true,urlcolor=moss,linkcolor=moss}
\newcommand{\playground}[3]{%
  \par\vspace{0.4ex}%
  {\footnotesize\color{ash}$\rightarrow$~\href{#1}{\bfseries #2} — #3\par}}

\newcommand{\кейс}[1]{{\color{brick}\small\textbf{#1}}}
\newcommand{\числ}[1]{{\color{navy}\bfseries #1}}

%
% Здесь только оформление и макросы — ни \title, ни текста слайдов, чтобы
% пять презентаций курса выглядели одинаково и правились в одном месте.
%
\usepackage{fontspec}
\usepackage{polyglossia}
\setmainlanguage{russian}
\setotherlanguage{english}
\setmainfont{Calibri}
\setsansfont{Calibri}
\setmonofont{Consolas}[Scale=0.92]

\usepackage{graphicx}
\usepackage{booktabs}
\usepackage{array}

% ---------------------------------------------------------------- оформление
\definecolor{ink}{HTML}{1B2433}
\definecolor{navy}{HTML}{2B4C6F}
\definecolor{brick}{HTML}{B03A2B}
\definecolor{moss}{HTML}{3E7D5A}
\definecolor{plum}{HTML}{6B5B95}
\definecolor{ash}{HTML}{8B9299}
\definecolor{paper}{HTML}{ECEFF3}

\usetheme{default}
\usecolortheme{default}
\setbeamercolor{structure}{fg=navy}
\setbeamercolor{frametitle}{fg=navy}
\setbeamercolor{normal text}{fg=ink}
\setbeamercolor{itemize item}{fg=navy}
\setbeamercolor{itemize subitem}{fg=ash}
\setbeamerfont{frametitle}{size=\Large,series=\bfseries}
\setbeamerfont{framesubtitle}{size=\small,series=\mdseries}
\setbeamercolor{framesubtitle}{fg=ash}
\setbeamertemplate{navigation symbols}{}
\setbeamertemplate{itemize item}{\raisebox{0.15ex}{\tiny$\bullet$}}
\setbeamertemplate{itemize subitem}{\raisebox{0.15ex}{\tiny--}}
\setbeamersize{text margin left=1.1cm, text margin right=1.1cm}

\setbeamertemplate{footline}{%
  \hbox{\begin{beamercolorbox}[wd=\paperwidth,ht=2.6ex,dp=1.4ex,leftskip=1.1cm,%
        rightskip=1.1cm]{}%
    {\color{ash}\scriptsize Основы машинного обучения \quad·\quad
     Цифровая лингвистика и локализация}\hfill
    {\color{ash}\scriptsize\insertframenumber}%
  \end{beamercolorbox}}}

% Врезка: связь слайда с лабораторной или с работой переводчика.
% Ширина считается точно: colorbox добавляет \fboxsep с каждой стороны,
% поэтому parbox должен быть на 2\fboxsep уже текстового поля — иначе
% на каждом слайде вылезает overfull hbox.
\newcommand{\врезка}[1]{%
  \par\vspace{0.5ex}\noindent
  \colorbox{paper}{%
    \parbox{\dimexpr\textwidth-2\fboxsep\relax}{%
      \vspace{0.7ex}\small #1\vspace{0.7ex}}}\par}

% Колонка таблицы заданной доли ширины, без растягивания пробелов.
% Объявлять её нужно здесь: внутри frame beamer пересканирует тело кадра,
% и #1 в определении вызывает «Illegal parameter number».
\newcolumntype{Л}[1]{>{\raggedright\arraybackslash}p{#1\textwidth}}

% Ссылка на интерактивную страницу к слайду. Playground'ы считают те же
% числа на том же корпусе, что и лабораторные, поэтому на них можно
% ссылаться прямо с лекции, не оговаривая расхождений.
\hypersetup{colorlinks=true,urlcolor=moss,linkcolor=moss}
\newcommand{\playground}[3]{%
  \par\vspace{0.4ex}%
  {\footnotesize\color{ash}$\rightarrow$~\href{#1}{\bfseries #2} — #3\par}}

\newcommand{\кейс}[1]{{\color{brick}\small\textbf{#1}}}
\newcommand{\числ}[1]{{\color{navy}\bfseries #1}}

%
% Здесь только оформление и макросы — ни \title, ни текста слайдов, чтобы
% пять презентаций курса выглядели одинаково и правились в одном месте.
%
\usepackage{fontspec}
\usepackage{polyglossia}
\setmainlanguage{russian}
\setotherlanguage{english}
\setmainfont{Calibri}
\setsansfont{Calibri}
\setmonofont{Consolas}[Scale=0.92]

\usepackage{graphicx}
\usepackage{booktabs}
\usepackage{array}

% ---------------------------------------------------------------- оформление
\definecolor{ink}{HTML}{1B2433}
\definecolor{navy}{HTML}{2B4C6F}
\definecolor{brick}{HTML}{B03A2B}
\definecolor{moss}{HTML}{3E7D5A}
\definecolor{plum}{HTML}{6B5B95}
\definecolor{ash}{HTML}{8B9299}
\definecolor{paper}{HTML}{ECEFF3}

\usetheme{default}
\usecolortheme{default}
\setbeamercolor{structure}{fg=navy}
\setbeamercolor{frametitle}{fg=navy}
\setbeamercolor{normal text}{fg=ink}
\setbeamercolor{itemize item}{fg=navy}
\setbeamercolor{itemize subitem}{fg=ash}
\setbeamerfont{frametitle}{size=\Large,series=\bfseries}
\setbeamerfont{framesubtitle}{size=\small,series=\mdseries}
\setbeamercolor{framesubtitle}{fg=ash}
\setbeamertemplate{navigation symbols}{}
\setbeamertemplate{itemize item}{\raisebox{0.15ex}{\tiny$\bullet$}}
\setbeamertemplate{itemize subitem}{\raisebox{0.15ex}{\tiny--}}
\setbeamersize{text margin left=1.1cm, text margin right=1.1cm}

\setbeamertemplate{footline}{%
  \hbox{\begin{beamercolorbox}[wd=\paperwidth,ht=2.6ex,dp=1.4ex,leftskip=1.1cm,%
        rightskip=1.1cm]{}%
    {\color{ash}\scriptsize Основы машинного обучения \quad·\quad
     Цифровая лингвистика и локализация}\hfill
    {\color{ash}\scriptsize\insertframenumber}%
  \end{beamercolorbox}}}

% Врезка: связь слайда с лабораторной или с работой переводчика.
% Ширина считается точно: colorbox добавляет \fboxsep с каждой стороны,
% поэтому parbox должен быть на 2\fboxsep уже текстового поля — иначе
% на каждом слайде вылезает overfull hbox.
\newcommand{\врезка}[1]{%
  \par\vspace{0.5ex}\noindent
  \colorbox{paper}{%
    \parbox{\dimexpr\textwidth-2\fboxsep\relax}{%
      \vspace{0.7ex}\small #1\vspace{0.7ex}}}\par}

% Колонка таблицы заданной доли ширины, без растягивания пробелов.
% Объявлять её нужно здесь: внутри frame beamer пересканирует тело кадра,
% и #1 в определении вызывает «Illegal parameter number».
\newcolumntype{Л}[1]{>{\raggedright\arraybackslash}p{#1\textwidth}}

% Ссылка на интерактивную страницу к слайду. Playground'ы считают те же
% числа на том же корпусе, что и лабораторные, поэтому на них можно
% ссылаться прямо с лекции, не оговаривая расхождений.
\hypersetup{colorlinks=true,urlcolor=moss,linkcolor=moss}
\newcommand{\playground}[3]{%
  \par\vspace{0.4ex}%
  {\footnotesize\color{ash}$\rightarrow$~\href{#1}{\bfseries #2} — #3\par}}

\newcommand{\кейс}[1]{{\color{brick}\small\textbf{#1}}}
\newcommand{\числ}[1]{{\color{navy}\bfseries #1}}


\title{Лабораторная работа № 4}
\subtitle{Генерация текста}
\date{}
\author{}

\begin{document}

% =====================================================================
\begin{frame}[plain]
  \vspace{1.2cm}
  {\color{brick}\small\bfseries ЛАБОРАТОРНАЯ РАБОТА № 4}

  \vspace{0.7cm}
  {\color{navy}\fontsize{30}{34}\selectfont\bfseries Генерация текста\par}

  \vspace{0.5cm}
  {\color{ink}\large Чем «складно» отличается от «выучено наизусть»\par}

  \vspace{1.0cm}
  {\color{ash} Работа — демонстрация на одно занятие, а не расчётная.
   Её смысл в том, чтобы увидеть устройство языковой модели целиком и
   понять, почему сегодня так уже не делают.\par}
\end{frame}

% =====================================================================
\begin{frame}{Языковая модель — это одно распределение}
  Вся языковая модель, от посимвольной цепочки до GPT, делает ровно одно:
  по уже написанному тексту выдаёт {\bfseries распределение вероятностей
  следующего элемента}.

  \vspace{0.9em}
  \begin{columns}[T,onlytextwidth]
    \column{0.47\textwidth}
      {\color{navy}\bfseries Что видит модель}\\[0.4em]
      {\small Контекст «Сохран» — и таблица по всем символам корпуса:
      «и» — 0,76, «е» — 0,12, «я» — 0,11, остальные почти ноль.}
    \column{0.47\textwidth}
      {\color{brick}\bfseries Что делает человек}\\[0.4em]
      {\small Выбирает, как из этой таблицы достать один символ. Это
      отдельное решение, к обучению модели отношения не имеющее — и
      именно оно решает, каким получится текст.}
  \end{columns}

  \vspace{1.0em}
  \врезка{\textbf{Генерация — это цикл.} Взять распределение, вытащить
  символ, дописать его к контексту, повторить. Больше в генерации текста
  ничего нет.}
\end{frame}

% =====================================================================
\begin{frame}{Почему в этом курсе модель посимвольная}
  \begin{columns}[T,onlytextwidth]
    \column{0.47\textwidth}
      {\color{navy}\bfseries Практические причины}\\[0.4em]
      {\small Корпус курса — 160 сегментов, около 7 тысяч символов. Для
      модели по словам это ничтожно мало.\par
      \vspace{0.6em}
      По символам алфавит маленький, данных хватает, и всё считается в
      браузере за секунды.}
    \column{0.47\textwidth}
      {\color{brick}\bfseries Содержательная причина}\\[0.4em]
      {\small Посимвольная модель прозрачна насквозь: она предсказывает
      следующий символ по предыдущим \числ{n} — и больше ничего.\par
      \vspace{0.6em}
      Всё, что дальше будет сказано про большие модели, проверяется прямо
      на ней.}
  \end{columns}

  \vspace{1.0em}
  \врезка{\textbf{Сглаживание обязательно.} Если сочетание символов не
  встретилось в обучении, его вероятность будет нулевой — и перплексия на
  любом новом тексте станет бесконечной. Сглаживание раздаёт невидённым
  продолжениям маленькую вероятность.}
\end{frame}

% =====================================================================
\begin{frame}{Температура: один параметр, три режима}
  \begin{columns}[T,onlytextwidth]
    \column{0.31\textwidth}
      {\color{navy}\bfseries T $<$ 1}\\[0.4em]
      {\small Распределение заостряется. При T $\to$ 0 модель всегда
      берёт самый вероятный символ.\par
      \vspace{0.5em}
      Текст предсказуемый, часто зацикливается.}
    \column{0.31\textwidth}
      {\color{ash}\bfseries T = 1}\\[0.4em]
      {\small Распределение как есть. Символы берутся ровно с теми
      вероятностями, которые выучила модель.}
    \column{0.31\textwidth}
      {\color{brick}\bfseries T $>$ 1}\\[0.4em]
      {\small Распределение сглаживается, редкие символы получают шанс.\par
      \vspace{0.5em}
      Текст разнообразнее и быстро перестаёт быть текстом.}
  \end{columns}

  \vspace{0.7em}
  \врезка{\textbf{Температура ничего не знает о смысле:} она меняет форму
  распределения, и только. «Творческий» режим большой модели — тот же
  множитель на тех же числах.}

  \playground{https://konkin-nikita.ru/gd_learn/lm/}{языковая модель и температура}%
  {ползунок температуры и распределение следующего символа рядом.}
\end{frame}

% =====================================================================
\begin{frame}{top-k и top-p: обрезать хвост}
  Температура растягивает всё распределение целиком. Обычно хочется
  другого: выбросить заведомо негодные продолжения, не трогая остальные.

  \vspace{0.9em}
  \begin{columns}[T,onlytextwidth]
    \column{0.47\textwidth}
      {\color{navy}\bfseries top-k}\\[0.4em]
      {\small Оставить \числ{k} самых вероятных символов, остальным дать
      ноль.\par
      \vspace{0.5em}
      Число выживших фиксировано — и когда модель уверена, и когда нет.}
    \column{0.47\textwidth}
      {\color{navy}\bfseries top-p}\\[0.4em]
      {\small Оставить столько символов, сколько нужно, чтобы их суммарная
      вероятность дошла до \числ{p}.\par
      \vspace{0.5em}
      Модель уверена — выживает один символ; сомневается — выживают
      десятки.}
  \end{columns}

  \vspace{0.7em}
  \врезка{Порядок применения важен: сначала температура, потом обрезка.
  Иначе обрезка уберёт то, что температура только собиралась поднять.}
\end{frame}

% =====================================================================
\begin{frame}{Перплексия: чем её меряют}
  \begin{columns}[T,onlytextwidth]
    \column{0.47\textwidth}
      {\small
      Перплексия — это «сколько вариантов модель в среднем перебирает на
      каждом символе».\par
      \vspace{0.7em}
      Перплексия 2 — модель колеблется между двумя символами.
      Перплексия 30 — практически не знает, что будет дальше.\par
      \vspace{0.7em}
      Меряется на тексте, {\bfseries которого модель не видела}. На
      обучающем тексте она меряет память, а не язык.}
    \column{0.47\textwidth}
      \врезка{\textbf{Отсюда — разделение корпуса.} Часть сегментов
      откладывается и в обучении не участвует. Всё, что говорится о
      качестве модели, говорится по отложенной части.\par
      \vspace{0.5em}
      Это та же дисциплина, что кросс-валидация в работе № 1.}
  \end{columns}
\end{frame}

% =====================================================================
\begin{frame}{Кривая переобучения}
  \framesubtitle{Символьная модель на корпусе курса, порядок от 1 до 8}

  \vspace{0.2em}
  {\small
  \begin{tabular}{@{}r@{\hspace{2.2em}}r@{\hspace{2.2em}}r@{\hspace{2.2em}}r@{}}
    \toprule
    порядок & перплексия на обучении & на отложенном & списано дословно \\
    \midrule
    1 & 13,37 & 16,14 & 0~\% \\
    2 & 5,81 & 11,09 & 0~\% \\
    \textbf{3} & \textbf{2,70} & {\color{moss}\bfseries 10,36} & \textbf{7~\%} \\
    4 & 1,80 & 11,67 & 66~\% \\
    5 & 1,51 & 13,37 & 84~\% \\
    6 & 1,42 & 14,90 & 100~\% \\
    8 & 1,39 & 17,13 & 100~\% \\
    \bottomrule
  \end{tabular}}

  \vspace{0.7em}
  {\small На обучении перплексия падает всё время. На отложенном тексте у
  неё {\color{brick}\bfseries минимум в середине} — на порядке 3. Дальше
  модель не учится языку, а запоминает корпус.}
\end{frame}

% =====================================================================
\begin{frame}{«Складно» и «выучено» — это разные вещи}
  \framesubtitle{Что модель выдаёт при разных порядках}

  \vspace{0.3em}
  {\footnotesize
  {\color{ash}порядок 1, списано 0~\%}\par
  \texttt{ереалитвеныевомачёть пстзUTTошетатм пмон / Прагранвыте ему}\par
  \vspace{0.6em}
  {\color{moss}порядок 3, списано 7~\% — лучший по перплексии}\par
  \texttt{Времено / Очисляется поное влияются настро ном реждый делю.}\par
  \vspace{0.6em}
  {\color{brick}порядок 6, списано 100~\%}\par
  \texttt{Временные файлы удаляются запрос должен в течение сохранен}}

  \vspace{0.9em}
  \врезка{\textbf{На глаз лучший — третий текст.} Он целиком составлен из
  кусков обучающего корпуса: ни одного своего сочетания символов. Модель,
  которая лучше всех обобщает, выглядит хуже той, которая просто помнит.}

  \playground{https://konkin-nikita.ru/gd_learn/lm/}{языковая модель и температура}%
  {списанные куски подсвечиваются прямо в сгенерированном тексте.}
\end{frame}

% =====================================================================
\begin{frame}{Почему это важно за пределами учебной модели}
  \begin{columns}[T,onlytextwidth]
    \column{0.47\textwidth}
      {\color{navy}\bfseries Механизм тот же}\\[0.4em]
      {\small Большая модель отличается размером и архитектурой, но
      «правдоподобно» у неё точно так же не означает «верно» и не
      означает «ново».\par
      \vspace{0.6em}
      Беглость текста — это свойство распределения, а не признак
      понимания.}
    \column{0.47\textwidth}
      {\color{brick}\bfseries Что из этого следует в работе}\\[0.4em]
      {\small Гладкий текст от языковой модели нельзя принимать как
      готовый перевод или готовую документацию.\par
      \vspace{0.6em}
      Вопрос «не воспроизведён ли дословно чужой текст» — не
      теоретический: на этой модели он измеряется прямо.}
  \end{columns}

  \vspace{1.0em}
  \врезка{Доля дословно списанного — это ровно та величина, вокруг
  которой идут все споры об авторском праве и больших моделях. Здесь её
  видно на корпусе в семь тысяч символов.}
\end{frame}

% =====================================================================
\begin{frame}{Что делает ноутбук работы}
  \begin{columns}[T,onlytextwidth]
    \column{0.47\textwidth}
      {\color{navy}\bfseries Что там происходит}\\[0.4em]
      {\small Посимвольная рекуррентная сеть обучается на текстах
      Шекспира — это учебник TensorFlow, встроенный в работу как есть.\par
      \vspace{0.6em}
      Десять эпох, несколько минут на GPU в Colab. Перебора
      гиперпараметров нет намеренно.}
    \column{0.47\textwidth}
      {\color{brick}\bfseries Почему так, а не иначе}\\[0.4em]
      {\small Посимвольная RNN — исторический этап, пройденный около 2015
      года. Настраивать её до приемлемого качества бессмысленно:
      приемлемого качества у этой архитектуры нет.\par
      \vspace{0.6em}
      Ценность в том, что её устройство целиком помещается в голову.}
  \end{columns}

  \vspace{0.9em}
  \врезка{Корпус там английский, и с корпусом курса результаты не
  сравнимы. Все числа этой лекции получены на русском корпусе курса
  символьной n-граммной моделью — её можно открыть в браузере и повторить.}
\end{frame}

% =====================================================================
\begin{frame}{Что должно быть в отчёте}
  \begin{columns}[T,onlytextwidth]
    \column{0.47\textwidth}
      {\color{moss}\bfseries Обязательно}\\[0.5em]
      {\small
      1. Что делает температура с распределением — своими словами, с
      примером текста при малой и большой.\par
      \vspace{0.4em}
      2. Где находится минимум перплексии на отложенном тексте и почему
      он не там, где текст красивее.\par
      \vspace{0.4em}
      3. Ответы на вопросы в конце ноутбука.}
    \column{0.47\textwidth}
      {\color{brick}\bfseries Не принимается}\\[0.5em]
      {\small «Модель сгенерировала осмысленный текст».\par
      \vspace{0.6em}
      Нужно сказать, на каком основании текст назван осмысленным, и
      проверить, не списан ли он. Для этого и нужна доля дословных
      совпадений.}
  \end{columns}
\end{frame}

% =====================================================================
\begin{frame}{Что унести из работы № 4}
  \begin{columns}[T,onlytextwidth]
    \column{0.47\textwidth}
      {\color{navy}\bfseries Про генерацию}\\[0.5em]
      {\small Языковая модель выдаёт распределение. Всё остальное —
      температура, top-k, top-p — решения о том, как из него выбирать.\par
      \vspace{0.6em}
      Качество текста на глаз и качество модели по перплексии — разные
      величины, и растут они в разные стороны.}
    \column{0.47\textwidth}
      {\color{brick}\bfseries Про профессию}\\[0.5em]
      {\small Беглый текст — слабейшее из возможных свидетельств качества.
      Проверять приходится то же, что и в работе № 3: фактическую
      точность, целостность, происхождение.}
  \end{columns}

  \vspace{1.2em}
  \playground{https://konkin-nikita.ru/gd_learn/lm/}{языковая модель и температура}%
  {порядок, температура, перплексия и доля списанного — в браузере.}
\end{frame}

\end{document}
